% Section 12.4: The Floating Cutoff.
% Source: physical PDF pp. 548--551 (printed pp. 525--528).
% Equations: (12.4.1)--(12.4.10). Figures/tables: none.
% Footnotes: four, including the section-title note.

\subsection[The Floating Cutoff]{The Floating Cutoff\texorpdfstring{\protect\footnotemark}{}}
\footnotetext{This section lies somewhat out of the book's main line of
development, and may be omitted in a first reading.}

Before closing this chapter, it is worth commenting on the relation between
conventional renormalization theory and an approach pioneered by
Wilson~\hyperref[ref:12.18]{[18]}. In Wilson's method one imposes a
``floating'' finite ultraviolet cutoff (either sharp or smooth) at momenta with
components of order $\Lambda$, and instead of taking $\Lambda\to\infty$, one
requires that the bare constants of the theory (those appearing in the
Lagrangian) depend on $\Lambda$ in such a way that all observable quantities
are $\Lambda$-independent.

It is convenient to work with dimensionless parameters. If a bare coupling or
mass parameter $g_i(\Lambda)$ has dimensionality $[\text{mass}]^{\Delta_i}$,
we define the corresponding dimensionless parameter $\mathcal G_i$ by
\begin{equation}
\mathcal G_i(\Lambda)\equiv\Lambda^{-\Delta_i}g_i(\Lambda).
\tag{12.4.1}\label{eq:12.4.1}
\end{equation}
Ordinary dimensional analysis tells us that the value of $\mathcal G_i$ at one
value $\Lambda'$ of the cutoff can be expressed as a function of the values of
the $\mathcal G_j$ at another value $\Lambda$ of the cutoff, and the
dimensionless ratio $\Lambda'/\Lambda$:
\begin{equation}
\mathcal G_i(\Lambda')=F_i\bigl(\mathcal G(\Lambda),\Lambda'/\Lambda\bigr).
\tag{12.4.2}\label{eq:12.4.2}
\end{equation}
No dimensional parameters other than $\Lambda'$ and $\Lambda$ can appear in
$F$, because no ultraviolet or infrared divergences can enter here; the
difference between the constants at $\Lambda$ and at $\Lambda'$ arises from
diagrams whose internal lines are restricted to have momenta between
$\Lambda$ and $\Lambda'$. Differentiating Eq.~\eqref{eq:12.4.2} with respect
to $\Lambda'$ and then setting $\Lambda'$ equal to $\Lambda$ yields a
differential equation for $\mathcal G_i$:
\begin{equation}
\Lambda\frac{\dd}{\dd\Lambda}\mathcal G_i(\Lambda)
=\beta_i\bigl(\mathcal G(\Lambda)\bigr),
\tag{12.4.3}\label{eq:12.4.3}
\end{equation}
where $\beta_i(\mathcal G)\equiv
[\partial F_i(\mathcal G,z)/\partial z]_{z=1}$. The functions
$\beta_i(\mathcal G)$ may be calculated for small couplings in perturbation
theory. This is Wilson's version of the ``renormalization group'' equation,
which will be discussed in somewhat different terms in Volume II.

The Lagrangian for any finite value of the cutoff defines an effective field
theory, in which instead of (or in addition to) integrating out ``heavy''
particles like the electron in the work of Euler \emph{et al.}, one integrates
out \emph{all} particles with momenta greater than $\Lambda$. Even if one
starts with a theory with a finite number of coupling parameters
$\mathcal G_i^0$ at some cutoff $\Lambda_0$, at any other value of the cutoff
the differential equation \eqref{eq:12.4.3} will generally yield non-zero
values for all couplings allowed by symmetry principles.\footnote{The only
known exceptions to this rule are in theories based on supersymmetry
\hyperref[ref:12.4]{[4]}.}

We now distinguish between the renormalizable and unrenormalizable couplings,
labelled $\mathcal G_a$ and $\mathcal G_n$, respectively, with $a$ running over
the finite number $N$ of couplings (including masses) for which
$\Delta_a\geq 0$, and $n$ running over the infinite number of couplings with
dimensionalities $\Delta_n<0$. We want to show that if the couplings
$\mathcal G_a(\Lambda_0)$ and $\mathcal G_n(\Lambda_0)$ at some initial cutoff
value $\Lambda_0$ lie on a generic $N$-dimensional initial surface
$\mathcal S_0$, then (with some qualifications) for $\Lambda\ll\Lambda_0$ they
will approach a fixed surface $\mathcal S$ that is independent of both
$\Lambda_0$ and the initial surface.\footnote{This theorem is due to
Polchinski~\hyperref[ref:12.19]{[19]}. What follows here is a shortened and less
rigorous version. (In Polchinski's proof, the initial surface is taken to be
that with all non-renormalizable couplings vanishing. As we shall see here,
the couplings approach the same fixed surface for generic initial surfaces.)}
This fixed surface is stable, in the sense that from any point on the surface,
the trajectory generated by Eq.~\eqref{eq:12.4.3} stays on the surface. Such a
stable surface defines a finite-parameter set of theories whose physical
content is cutoff-independent, which as argued in the previous section, is the
essential property of renormalizable theories. Furthermore, this construction
shows that a generic theory defined with cutoff $\Lambda_0$ will look for
$\Lambda\ll\Lambda_0$ like a renormalizable theory.\footnote{Of course some
theories have symmetries and a field content that do not allow any
renormalizable interactions. This is the case for theories containing only
fermion fields, or only the gravitational field. Such theories for
$\Lambda\ll\Lambda_0$ look like free-field theories.}

To prove these results, consider any small perturbation $\delta\mathcal
G_i(\Lambda)$ in the values of the $\mathcal G_i(\Lambda)$ satisfying
Eq.~\eqref{eq:12.4.3}. It will satisfy the differential equation
\begin{equation}
\Lambda\frac{\dd}{\dd\Lambda}\delta\mathcal G_i(\Lambda)
=\sum_j M_{ij}\bigl(\mathcal G(\Lambda)\bigr)\delta\mathcal G_j(\Lambda),
\tag{12.4.4}\label{eq:12.4.4}
\end{equation}
where
\begin{equation}
M_{ij}(\mathcal G)\equiv\frac{\partial}{\partial\mathcal G_j}\beta_i(\mathcal G).
\tag{12.4.5}\label{eq:12.4.5}
\end{equation}
This equation couples the renormalizable and unrenormalizable couplings, making
it difficult to see the difference in their behavior. To decouple them, we
introduce the linear combinations
\begin{equation}
\xi_n\equiv\delta\mathcal G_n-
\sum_{ab}\frac{\partial\mathcal G_n}{\partial\mathcal G_a^0}
\left(\frac{\partial\mathcal G}{\partial\mathcal G^0}\right)^{-1}_{ab}
\delta\mathcal G_b,
\tag{12.4.6}\label{eq:12.4.6}
\end{equation}
where $\mathcal G_a^0$ are the values of the renormalizable couplings at cutoff
$\Lambda_0$, which we shall use as coordinates for the initial surface, and
$\mathcal G_n$ are the values of the non-renormalizable couplings at cutoff
$\Lambda$ derived from the differential equation \eqref{eq:12.4.3}, with
initial value for $\Lambda_0$ at the point on the initial surface with
coordinates $\mathcal G_a^0$. To calculate the derivative of $\xi_n$ with
respect to $\Lambda$, we note that the derivatives
$\partial\mathcal G_i/\partial\mathcal G_a^0$ satisfy the same differential
equation \eqref{eq:12.4.4} as the $\delta\mathcal G_i$. It is an elementary
exercise then to show that
\begin{equation}
\Lambda\frac{\dd}{\dd\Lambda}\xi_n=\sum_m N_{nm}\xi_m,
\tag{12.4.7}\label{eq:12.4.7}
\end{equation}
where
\begin{equation}
N_{nm}\equiv M_{nm}-\sum_{ab}
\frac{\partial\mathcal G_n}{\partial\mathcal G_a^0}
\left(\frac{\partial\mathcal G}{\partial\mathcal G^0}\right)^{-1}_{ab}M_{bm}.
\tag{12.4.8}\label{eq:12.4.8}
\end{equation}
Now we must estimate the elements of $N_{nm}$. For a free-field theory no
cutoff is needed, so for very small coupling all bare parameters $g_i(\Lambda)$
become $\Lambda$-independent. Hence for small couplings the dimensionless
parameters $\mathcal G_i$ simply scale as $\Lambda^{-\Delta_i}$, and the matrix
$M_{ij}$ is given by
\begin{equation}
M_{ij}\approx-\Delta_i\delta_{ij}.
\tag{12.4.9}\label{eq:12.4.9}
\end{equation}
It follows that the matrix $N_{nm}$ is given approximately by
$-\Delta_n\delta_{nm}$. The defining characteristic of the
non-renormalizable couplings is that $\Delta_n<0$, so Eq.~\eqref{eq:12.4.7}
tells us that, at least for couplings in some finite range, where $N_{nm}$ is
positive-definite, the $\xi_n$ decay for $\Lambda\ll\Lambda_0$ like positive
powers of $\Lambda/\Lambda_0$. In this limit, then, the perturbations are
related by
\begin{equation}
\delta\mathcal G_n=\sum_{ab}
\frac{\partial\mathcal G_n}{\partial\mathcal G_a^0}
\left(\frac{\partial\mathcal G}{\partial\mathcal G^0}\right)^{-1}_{ab}
\delta\mathcal G_b.
\tag{12.4.10}\label{eq:12.4.10}
\end{equation}
In particular, if we make a small change in the initial surface $\mathcal S_0$
and/or the starting point on that surface and/or the initial cutoff
$\Lambda_0$, such that the perturbations $\delta\mathcal G_a$ in the
renormalizable couplings vanish at some cutoff $\Lambda\ll\Lambda_0$, then the
perturbations $\delta\mathcal G_n$ in all the other couplings at cutoff
$\Lambda$ also vanish. Thus the non-renormalizable couplings
$\mathcal G_n(\Lambda)$ for $\Lambda\ll\Lambda_0$ can depend only on the
renormalizable couplings $\mathcal G_a(\Lambda)$, not separately on the initial
surface or the starting point on that surface or the initial cutoff
$\Lambda_0$. At cutoff $\Lambda\ll\Lambda_0$ all the couplings therefore
approach an $N$-dimensional surface $\mathcal S$, with coordinates
$\mathcal G_a(\Lambda)$, which is independent of both the initial surface and
of $\Lambda_0$. Note that the non-renormalizable couplings $\mathcal G_n$ are
not generally small on $\mathcal S$; the important point is that they become
functions of the renormalizable couplings. Changes in $\Lambda$ with $\Lambda$
remaining much less than $\Lambda_0$ will change the couplings, but the
couplings will remain close to $\mathcal S$ (at least as long as the couplings
do not become so large that $N_{nm}$ is no longer a positive-definite matrix).
Hence $\mathcal S$ is a stable surface, as was to be proved.

We have seen that all physical quantities may be expressed in terms of
$\Lambda$ and the $\mathcal G_n(\Lambda)$, and are $\Lambda$-independent. This
is true in particular of the $N$ conventional renormalized couplings and
masses, like $e$ and $m_e$ in quantum electrodynamics. But we can then invert
this relation, and express the $\mathcal G_n(\Lambda)$ in terms of the
conventional parameters and $\Lambda$. In this way we can justify the usual
renormalization program; all physical quantities are expressed in a
cutoff-independent way in terms of the conventional renormalized couplings and
masses.

The Wilson approach has some advantages in practice. One does not have to
worry about subintegrations and overlapping divergences; the momentum cutoff
applies to all internal lines. Also, some of the non-renormalization theorems
of supersymmetry theory, which tell us that certain couplings are not affected
by radiative corrections, work only for the cutoff-dependent bare
couplings~\hyperref[ref:12.20]{[20]}.

On the other hand, there are disadvantages to the Wilson approach. One must
give up the special simplicities of working with renormalizable theories like
quantum electrodynamics; once one starts integrating out particles with
momenta above some scale $\Lambda$, the resulting effective field theory will
contain all Lorentz- and gauge-invariant interactions, with $\Lambda$-dependent
couplings. (Nevertheless, in physical processes at energies $E\ll\Lambda$,
the dominant couplings will still be the renormalizable ones.) Also, the
cutoff generally destroys \emph{manifest} gauge invariance, and either manifest
Lorentz invariance or unitarity. None of this is a problem in condensed matter
physics, the original context of Wilson's approach, because no one would
expect a realistic condensed matter theory to be strictly renormalizable, and
there are no fundamental physical principles that are necessarily violated by
a cutoff. In fact, in crystals there is a cutoff on phonon momenta, provided by
the inverse lattice spacing.

At bottom, the difference between the conventional and the Wilson approach is
one of mathematical convenience rather than of physical interpretation.
Indeed, conventional renormalization already provides a sort of adjustable
cutoff; when we express our answer in terms of coupling constants that are
defined as the values of physical amplitudes at some momenta of order $\mu$
(as for the scalar field theory discussed in the previous section), the
cancellations that make integrals converge begin to operate at virtual momenta
of order $\mu$. Conversely, the $\Lambda$-dependent coupling constants of the
Wilson approach must ultimately be expressed in terms of observable masses and
charges, and when this is done the results are, of course, the same as those
obtained by conventional means.

